\begin{abstract}
Test-time adaptive model merging has emerged as a promising paradigm for combining specialized task-expert neural networks without retraining, utilizing unlabeled test-time streams. However, existing state-of-the-art frameworks introduce substantial architectural complexity, relying on overparameterized normalizing flows (e.g., 2.6M parameters in FoldMerge) or high-dimensional layer-wise adapters (e.g., in SyMerge and AdaMerging). In this work, guided by Occam's razor, we present a critical, deconstructive audit of test-time adaptive model merging, exploring the fundamental boundaries of parameter frugality. We show that for extremely low-parameter regimes (like our 8 global task-scalars), weight-space optimization alone is too mathematically constrained to resolve parameter conflicts, making classification head adaptation the primary driver of performance. Conversely, higher-capacity layer-wise methods achieve genuine weight-space gains by utilizing more degrees of freedom. To isolate these architectural boundaries, we introduce Barycentric Proximity-Anchored Merging (BPAM), a simple, low-parameter framework. BPAM optimizes only a single global scalar coefficient per task (exactly $K=8$ parameters total for an 8-task benchmark), restricting coefficients to a convex barycentric simplex and stabilizing optimization via a closed-form Mean-Field Proximity Regularizer. Under strictly frozen classification heads, BPAM-Static (8 parameters on the whole image encoder) achieves 69.21\% average accuracy, matching linear Task Arithmetic, while restricting merging to a single projection layer (BPAM-Restricted) collapses performance to 51.38\%. Concurrently adapting classification heads (BPAM-Full) yields 75.22\% average accuracy. Our findings offer a transparent, minimalist audit of test-time model merging, highlighting the crucial role of parameter constraints and clarifying the boundary where layer-wise degrees of freedom become necessary.
\end{abstract}
